\chapter{ฟิสิกส์ดาราศาสตร์และจักรวาลวิทยา (Astrophysics and Cosmology)}
\label{chap:astrophysics}
\index{ฟิสิกส์ดาราศาสตร์}
\index{จักรวาลวิทยา}
\index{หลุมดำ}
\index{สมการฟรีดมันน์}

\begin{chapterobjectives}
เมื่อศึกษาบทที่ 7 จบแล้ว นิสิตและผู้ศึกษาสามารถ:
\begin{enumerate}
    \item อธิบายสมดุลไฮโดรสแตติก (Hydrostatic Equilibrium) ในโครงสร้างภายในของดาวฤกษ์
    \item วิเคราะห์สถานะสุดท้ายของดาวฤกษ์ ได้แก่ ดาวแคระขาว ดาวนิวตรอน และขีดจำกัดจันทรเศขร (Chandrasekhar Limit)
    \item คำนวณรัศมีชวาร์ซชิลด์ (Schwarzschild Radius) และการยืดของเวลาเนื่องจากแรงโน้มถ่วง (Gravitational Time Dilation) ใกล้หลุมดำ
    \item อธิบายหลักการตรวจวัดคลื่นความโน้มถ่วง (Gravitational Waves) จากการรวมตัวของหลุมดำคู่
    \item ประยุกต์ใช้กฎของฮับเบิล-เลอแมทร์และสมการฟรีดมันน์ (Friedmann Equations) ในการอธิบายการขยายตัวของเอกภพ
    \item เขียนสคริปต์ Python จำลองการเปลี่ยนแปลงของสเกลแฟกเตอร์ $a(t)$ ในแบบจำลอง $\Lambda\text{CDM}$
\end{enumerate}
\end{chapterobjectives}

\section{สมดุลไฮโดรสแตติกและวิวัฒนาการของดาวฤกษ์}
\index{สมดุลไฮโดรสแตติก}
\index{ความดันดีเจเนอเรซี}

ดาวฤกษ์ดำรงสถานะเสถียรอยู่ได้ด้วยความสมดุลระหว่างแรงโน้มถ่วงที่พยายามดึงดูดมวลสารเข้าสู่ศูนย์กลาง กับแรงดันจากความร้อนและรังสีนิวเคลียร์ฟิวชันจากแกนกลาง เรียกว่า \textbf{สมดุลไฮโดรสแตติก (Hydrostatic Equilibrium)}:
\begin{equation}
\frac{dP}{dr} = -\frac{G M(r) \rho(r)}{r^2}
\end{equation}
โดยที่ $M(r) = \int_0^r 4\pi {r'}^2 \rho(r') dr'$ คือมวลสะสมภายในรัศมี $r$

\begin{figure}[H]
\centering
\begin{tikzpicture}[scale=1.0]
    % Stellar Core
    \draw[line width=1.5pt, darkNavy, fill=quantumGold!30] (0, 0) circle (2.2cm);
    \filldraw[red!80!black] (0, 0) circle (0.6cm);
    \node at (0, 0) [font=\tiny\bfseries\color{white}, align=center] {แกนฟิวชัน};
    
    % Gravity arrows (Inward)
    \foreach \angle in {30, 75, 120, 165, 210, 255, 300, 345} {
        \draw[-{Stealth[scale=1.2]}, blue!80!black, line width=1.4pt] 
            (\angle:3.2) -- (\angle:2.3);
    }
    \node at (2.6, 2.6) [font=\footnotesize\bfseries\color{blue!80!black}] {แรงโน้มถ่วงดึงดูด ($F_g$)};
    
    % Thermal Pressure arrows (Outward)
    \foreach \angle in {10, 55, 100, 145, 190, 235, 280, 325} {
        \draw[-{Stealth[scale=1.2]}, red!80!black, line width=1.4pt] 
            (\angle:0.8) -- (\angle:1.7);
    }
    \node at (-2.3, -2.4) [font=\footnotesize\bfseries\color{red!80!black}] {แรงดันความร้อน ($P$)};
\end{tikzpicture}
\caption{สภาวะสมดุลไฮโดรสแตติก (Hydrostatic Equilibrium) ของดาวฤกษ์}
\label{fig:hydrostatic_equilibrium}
\end{figure}

เมื่อดาวฤกษ์เผาผลาญเชื้อเพลิงนิวเคลียร์ที่แกนกลางจนหมด แรงดันความร้อนจะหยุดลง ดาวฤกษ์จะยุบตัวลงตามมวลตั้งต้น:
\begin{enumerate}
    \item \textbf{ดาวแคระขาว (White Dwarf):} พยุงตัวด้วย \textbf{ความดันดีเจเนอเรซีของอิเล็กตรอน (Electron Degeneracy Pressure)} ตามหลักการกีดกันของเพาลี มวลสูงสุดที่เป็นไปได้คือ \textbf{ขีดจำกัดจันทรเศขร (Chandrasekhar Limit)}:
    \begin{equation}
    M_{\text{Ch}} \approx 1.44 M_\odot
    \end{equation}
    \item \textbf{ดาวนิวตรอน (Neutron Star):} หากมวลเกิน $1.44 M_\odot$ อิเล็กตรอนจะรวมตัวกับโปรตอนกลายเป็นนิวตรอน พยุงตัวด้วยความดันดีเจเนอเรซีของนิวตรอน มวลสูงสุดอยู่ที่ขีดจำกัด TOV (Tolman-Oppenheimer-Volkoff Limit) $\approx 2.1 - 2.3 M_\odot$
    \item \textbf{หลุมดำ (Black Hole):} หากมวลสารตกค้างมากกว่าขีดจำกัด TOV ไม่มีแรงใดในเอกภพสามารถต้านทานแรงโน้มถ่วงได้ ดาวจะยุบตัวลงสู่ภาวะเอกฐาน (Singularity)
\end{enumerate}

\section{หลุมดำและเรขาคณิตชวาร์ซชิลด์ (Schwarzschild Black Holes)}
\index{หลุมดำ}
\index{รัศมีชวาร์ซชิลด์}
\index{ขอบฟ้าเหตุการณ์}

ตามทฤษฎีสัมพัทธภาพทั่วไปของไอน์สไตน์ คาร์ล ชวาร์ซชิลด์ (Karl Schwarzschild) ได้หาผลเฉลยแม่นตรงสำหรับปริภูมิ-เวลารอบมวลทรงกลมสถิตที่ไม่มีประจุและไม่หมุน:
\begin{equation}
ds^2 = -\left(1 - \frac{R_s}{r}\right) c^2 dt^2 + \left(1 - \frac{R_s}{r}\right)^{-1} dr^2 + r^2 (d\theta^2 + \sin^2\theta d\phi^2)
\end{equation}
โดยที่ $R_s$ คือ \textbf{รัศมีชวาร์ซชิลด์ (Schwarzschild Radius)} ซึ่งเป็นตำแหน่งของ \textbf{ขอบฟ้าเหตุการณ์ (Event Horizon)}:
\begin{equation}
R_s = \frac{2GM}{c^2}
\end{equation}
ไม่มีสิ่งใดแม้กระทั่งแสงสามารถหลุดรอดออกจากภายในขอบฟ้าเหตุการณ์ได้

การยืดของเวลาเนื่องจากแรงโน้มถ่วงสำหรับผู้สังเกตที่ระยะห่าง $r$ สัมพันธ์กับเวลาเฉพาะที่ระยะอนันต์คือ:
\begin{equation}
\Delta t(r) = \frac{\Delta t_\infty}{\sqrt{1 - \frac{R_s}{r}}}
\end{equation}
เมื่อ $r \to R_s$ เวลาจะหยุดเดินในสายตาของผู้สังเกตภายนอก ($\Delta t \to \infty$)

\section{จักรวาลวิทยาและสมการฟรีดมันน์ (Friedmann Cosmology)}
\index{กฎของฮับเบิล}
\index{สมการฟรีดมันน์}
\index{สสารมืด}
\index{พลังงานมืด}

เอ็ดวิน ฮับเบิล (Edwin Hubble) และจอร์จ เลอแมทร์ (Georges Lemaître) ค้นพบว่าดาราจักรส่วนใหญ่กำลังเคลื่อนที่ถอยห่างออกจากเราด้วยความเร็วแปรผันตรงกับระยะทาง:
\begin{equation}
v = H_0 d
\end{equation}
โดยที่ $H_0 \approx 70\text{ km}\cdot\text{s}^{-1}\cdot\text{Mpc}^{-1}$ คือค่าคงตัวของฮับเบิล (Hubble Constant)

ตามแบบจำลอง FLRW (Friedmann-Lemaître-Robertson-Walker Metric) พลวัตการขยายตัวของเอกภพถูกกำหนดด้วย \textbf{สเกลแฟกเตอร์ $a(t)$} ผ่านสมการฟรีดมันน์:
\begin{equation}
H^2(t) = \left(\frac{\dot{a}}{a}\right)^2 = \frac{8\pi G}{3} \rho - \frac{k c^2}{a^2} + \frac{\Lambda c^2}{3}
\end{equation}
และสมการความเร่ง:
\begin{equation}
\frac{\ddot{a}}{a} = -\frac{4\pi G}{3}\left(\rho + \frac{3P}{c^2}\right) + \frac{\Lambda c^2}{3}
\end{equation}
การสังเกตการณ์ซูเปอร์โนวา Type Ia ในปี ค.ศ. 1998 ยืนยันว่าเอกภพกำลังขยายตัวด้วยความเร่ง ($\ddot{a} > 0$) ขับเคลื่อนโดย \textbf{พลังงานมืด (Dark Energy / Cosmological Constant $\Lambda$)} ซึ่งครองสัดส่วนประมาณ 68\% ของพลังงานทั้งหมดในเอกภพ

\section{สถาปัตยกรรมการจำลองเสมือนจริง AR/XR: Spacetime Curvature \& Lensing}
\index{WebXR}
\index{Black Hole AR}

\begin{arbox}[การจำลองการบิดโค้งของปริภูมิเวลาและเลนส์โน้มถ่วงรอบหลุมดำ 3 มิติ]
\textbf{วัตถุประสงค์การเรียนรู้ผ่าน AR:}
\begin{enumerate}
    \item แสดงหลุมดำมวลยิ่งยวดพร้อมจานสะสมมวล (Accretion Disk) เรืองแสงและปรากฏการณ์เงาหลุมดำ (Black Hole Shadow)
    \item ใช้ GLSL Custom Shader จำลองการหักเหของเส้นทางแสง (Gravitational Lensing) จากดวงดาวพื้นหลังรอบขอบฟ้าเหตุการณ์
    \item ผู้เรียนสามารถโยนดาวเคราะห์หรืออนุภาคเสมือนเข้าสู่วงโคจร เพื่อสังเกตปรากฏการณ์เลื่อนไปข้างหน้าของจุดใกล้ศูนย์กลางที่สุด (Perihelion Precession) และการถูกฉีกขาดจากแรงไทดัล (Spaghettification)
\end{enumerate}
\end{arbox}

\section{การคำนวณเชิงตัวเลขด้วย Python: Friedmann Universe Expansion}
\index{Python}
\index{Cosmology Solver}

\begin{pythonbox}[friedmann\_cosmology.py: การจำลองการขยายตัวของเอกภพ Lambda-CDM]
import numpy as np
from scipy.integrate import solve_ivp

# พารามิเตอร์จักรวาลวิทยาปัจจุบัน (Planck Collaboration 2018)
H0 = 67.4       # km/s/Mpc
H0_per_s = H0 * 1000.0 / (3.08567758e22)  # แปลงเป็น 1/s
H0_per_Gyr = H0_per_s * (3.15576e16)       # แปลงเป็น 1/Gyr (~0.0689 Gyr^-1)

Omega_m = 0.315   # สสาร (สสารปกติ + สสารมืด)
Omega_r = 9.0e-5  # รังสี (โฟตอน + นิวทริโน)
Omega_L = 0.685   # พลังงานมืด (ค่าคงตัวจักรวาลวิทยา)

def friedmann_ode(t, a):
    """
    da/dt = H0 * sqrt( Omega_r/a^2 + Omega_m/a + Omega_L * a^2 )
    สำหรับเอกภพที่แบนราบทางเรขาคณิต (k = 0)
    """
    if a <= 0:
        return 0
    term = (Omega_r / a**2) + (Omega_m / a) + (Omega_L * a**2)
    return H0_per_Gyr * np.sqrt(max(term, 0))

# คำนวณย้อนหลังจากปัจจุบัน a=1 (t=13.8 Gyr)
t_span = (0.01, 25.0)  # เวลาตั้งแต่ 10 ล้านปี ถึง 25 พันล้านปี
a0 = [0.001]
t_eval = np.linspace(0.01, 25.0, 300)

sol = solve_ivp(friedmann_ode, t_span, a0, t_eval=t_eval)

# หาอายุเอกภพเมื่อ a = 1.0 (ปัจจุบัน)
curr_idx = np.argmin(np.abs(sol.y[0] - 1.0))
t_age = sol.t[curr_idx]

print("=== Friedmann Cosmological Simulation ===")
print(f"Current Hubble Constant H0: {H0} km/s/Mpc")
print(f"Calculated Age of Universe (a=1): {t_age:.2f} Billion Years (Gyr)")
\end{pythonbox}

\section{ตัวอย่างการแก้โจทย์เชิงวิเคราะห์ตามระเบียบวิธี P-S-C-T}

\begin{psctexample}[รัศมีชวาร์ซชิลด์และการยืดของเวลาใกล้หลุมดำใจกลางทางช้างเผือก]
\textbf{โจทย์ (Problem):} \\
หลุมดำมวลยิ่งยวด ซาจิททาเรียส เอ สตาร์ ($\text{Sgr A}^*$) ณ ใจกลางดาราจักรทางช้างเผือก มีมวลประมาณ $M = 4.15 \times 10^6 M_\odot$ (โดยมวลดวงอาทิตย์ $M_\odot \approx 1.989 \times 10^{30}\text{ kg}$)
\begin{enumerate}
    \item จงคำนวณรัศมีชวาร์ซชิลด์ $R_s$ ของหลุมดำนี้ในหน่วยกิโลเมตร
    \item หากยานสำรวจอวกาศลอยนิ่งอยู่ที่ระยะห่าง $r = 1.20 R_s$ เวลาที่ผ่านไป $1.0\text{ ชั่วโมง}$ บนยาน จะเทียบเท่ากับเวลากี่ชั่วโมงสำหรับผู้สังเกตที่อยู่ห่างไกลบนโลก?
\end{enumerate}

\textbf{ยุทธวิธี (Strategy):} \\
\begin{itemize}
    \item รัศมีชวาร์ซชิลด์: $R_s = \frac{2GM}{c^2}$
    \item การยืดของเวลาเนื่องจากแรงโน้มถ่วง: $\Delta t_\infty = \Delta t(r) \sqrt{1 - \frac{R_s}{r}}$ หรือ $\Delta t(r) = \frac{\Delta t_\infty}{\sqrt{1 - R_s/r}}$ \\
    (ระวัง: เวลาเฉพาะ $\Delta \tau = \Delta t(r)$ บนยานเดินช้ากว่าเวลาของผู้สังเกตไกลโพ้น $\Delta t_\infty = \frac{\Delta \tau}{\sqrt{1 - R_s/r}}$)
\end{itemize}

\textbf{การคำนวณ (Computation):} \\
\textit{1. คำนวณรัศมีชวาร์ซชิลด์:}
\begin{align}
M &= (4.15 \times 10^6) \times (1.989 \times 10^{30}\text{ kg}) \approx 8.254 \times 10^{36}\text{ kg} \\
R_s &= \frac{2 (6.674 \times 10^{-11}\text{ m}^3\text{kg}^{-1}\text{s}^{-2}) (8.254 \times 10^{36}\text{ kg})}{(2.998 \times 10^8\text{ m/s})^2} \\
&\approx \frac{1.1018 \times 10^{27}}{8.988 \times 10^{16}} \approx 1.226 \times 10^{10}\text{ m} \approx 1.226 \times 10^7\text{ km} \approx 12.26\text{ ล้านกิโลเมตร}
\end{align}
(สังเกต: รัศมีชวาร์ซชิลด์นี้มีขนาดประมาณ $0.082\text{ AU}$ ซึ่งเล็กกว่าวงโคจรของดาวพุธรอบดวงอาทิตย์!)

\textit{2. คำนวณการยืดของเวลา:}
ที่ระยะ $r = 1.20 R_s$:
\begin{equation}
\sqrt{1 - \frac{R_s}{r}} = \sqrt{1 - \frac{1}{1.20}} = \sqrt{1 - 0.8333} = \sqrt{0.1667} \approx 0.4082
\end{equation}
เวลาบนยาน $\Delta \tau = 1.0\text{ ชั่วโมง}$:
\begin{equation}
\Delta t_\infty = \frac{\Delta \tau}{\sqrt{1 - R_s/r}} = \frac{1.0\text{ hr}}{0.4082} \approx 2.45\text{ ชั่วโมง} \approx 2\text{ ชั่วโมง } 27\text{ นาที}
\end{equation}

\textbf{ข้อคิดทางฟิสิกส์ (Takeaway):} \\
สนามโน้มถ่วงมหาศาลทำให้โครงข่ายปริภูมิเวลาบิดโค้งอย่างรุนแรง เวลาบนยานที่อยู่ใกล้ขอบฟ้าเหตุการณ์เพียง $1.2 R_s$ จะเดินช้ากว่าเวลาของผู้สังเกตภายนอกมากกว่า 2.4 เท่า และหากยานขยับเข้าใกล้ขอบฟ้าเหตุการณ์ $r \to R_s$ เวลาภายนอกจะผ่านไปอย่างรวดเร็วจนเข้าสู่อนันต์
\qedexam
\end{psctexample}

\section{สรุปสาระสำคัญประจำบท}

\begin{table}[H]
\centering
\begin{tabularx}{\textwidth}{l Z Y}
\toprule
\textbf{หัวข้อ} & \textbf{สมการหลัก} & \textbf{ความหมายทางกายภาพ} \\
\midrule
สมดุลไฮโดรสแตติก & $\frac{dP}{dr} = -\frac{GM\rho}{r^2}$ & เสถียรภาพโครงสร้างของดาวฤกษ์ \\
ขีดจำกัดจันทรเศขร & $M_{\text{Ch}} \approx 1.44 M_\odot$ & มวลสูงสุดของดาวแคระขาว \\
รัศมีชวาร์ซชิลด์ & $R_s = \frac{2GM}{c^2}$ & ขอบเขตขอบฟ้าเหตุการณ์ของหลุมดำ \\
กฎของฮับเบิล & $v = H_0 d$ & เอกภพกำลังขยายตัวอย่างเป็นเอกรูป \\
สมการฟรีดมันน์ & $H^2 = \frac{8\pi G}{3}\rho - \frac{kc^2}{a^2} + \frac{\Lambda c^2}{3}$ & สมการพลวัตการขยายตัวของเอกภพ \\
\bottomrule
\end{tabularx}
\caption{สรุปความสัมพันธ์พื้นฐานในฟิสิกส์ดาราศาสตร์และจักรวาลวิทยา}
\label{tab:astrophysics_summary}
\end{table}

\section{แบบฝึกหัดท้ายบท}
\begin{enumerate}
    \item หากโลกถูกบีบอัดจนกลายเป็นหลุมดำ จงคำนวณรัศมีชวาร์ซชิลด์ของโลก กำหนดมวลโลก $M_E \approx 5.972 \times 10^{24}\text{ kg}$
    \item ดาราจักรแห่งหนึ่งอยู่ห่างจากโลก $100\text{ Mpc}$ จะมีความเร็วในการถอยห่างจากโลกเท่าใดตามกฎของฮับเบิล กำหนด $H_0 = 70\text{ km/s/Mpc}$
    \item จงอธิบายว่าเหตุใดความดันดีเจเนอเรซีของอิเล็กตรอนจึงไม่สามารถพยุงดาวที่มีมวลมากกว่า $1.44 M_\odot$ ได้ (เชื่อมโยงกับทฤษฎีสัมพัทธภาพพิเศษเมื่อความเร็วอิเล็กตรอนเข้าใกล้ $c$)
\end{enumerate}
